% ============================================================================
% Fichier  : partie6/P06_C21_auditologie_axiomes_piliers.tex
% Titre    : L'Auditologie : Objet, Axiomes et Piliers
% Auteur   : MINKA MI NGUIDJOÏ Thierry Emmanuel
% Version  : 1.0 -- novembre 2026
% Base   : ACSI C21, repris et développé pour INF3226
% ============================================================================
% Source unique pour le cadre : article A0 (hal-05701007, prépublication du
% 22/07/2026), version française : définitions de la mission et de
% l'auditologie, six axiomes ontologiques (énoncés repris, numérotation et
% intitulés d'A0), « ce qui n'est pas un axiome », quatre propositions,
% cinq piliers et leurs prédictions falsifiables, interactions entre
% piliers, science du design, engagement lakatossien (septième axiome,
% méthodologique), matrice de falsifiabilité, « ce que l'auditologie n'est
% pas », mini-cas LizbiZ (H ≈ 1,2 ; 5 × 5 ; usage stratégique de NEK ;
% 35 % de probabilité de régression à 18 mois ; H ≈ 0,6 six mois plus
% tard). Références classiques recopiées de la bibliographie d'A0.
% Environnement axiome : compteur global, utilisé ici seulement, de sorte
% que la numérotation 1 à 7 coïncide avec celle d'A0.
% ============================================================================

\chapter{L'Auditologie : Objet, Axiomes et Piliers}
\label{chap:auditologie}
\chaptermark{L'auditologie}

\epigraph{%
    « Le métier dit comment auditer.\\
    La science cherche pourquoi cela marche, ou échoue. »%
}{D'après l'article A0}

\niveauChapitre{Recherche}{%
    Maîtrise du Chapitre~\ref{chap:conclure} (mesures $\mu$, $U$, $H$,
    divergence) et du Chapitre~\ref{chap:decomposition} (gabarit
    canonique)%
}

\begin{boxobjectifs}
\begin{itemize}
    \item Définir la mission auditologique
          $\langle \mathcal{S}, \mathcal{A}, \mathcal{I}, \mathcal{G},
          \mathcal{T} \rangle$ et l'auditologie.
    \item Énoncer les six axiomes ontologiques, et ce qui n'est pas un
          axiome.
    \item Présenter les cinq piliers et leurs prédictions falsifiables.
    \item Situer la méthode du cours comme instanciation opérationnelle du
          programme, dans un cadre de science du design et de
          falsifiabilité lakatossienne.
\end{itemize}
\end{boxobjectifs}

\bigskip

% ============================================================================
\section*{Question directrice : pourquoi cette méthode~?}
% ============================================================================

Tout au long du cours, AUDIX CONSEIL a fait des choix qui n'allaient pas de
soi~: admettre la réponse \NEK{}, décomposer chaque domaine sous la borne de
sept, lire trois mesures au lieu d'un taux, comparer deux cycles,
considérer l'indépendance comme une condition de l'opposabilité. Le
Chapitre~\ref{chap:audit-auditologie} annonçait que ces choix découlaient
d'axiomes. Ce chapitre les présente pour eux-mêmes, tels que l'article A0
les formule~\cite{minka2026auditologie}.

\medskip
\textbf{Plan du chapitre.}
\begin{enumerate}
    \item L'objet~: la mission auditologique (section~\ref{sec:c30-objet}).
    \item Six axiomes ontologiques (section~\ref{sec:c30-axiomes}).
    \item Quatre conséquences, quatre pratiques (section~\ref{sec:c30-propositions}).
    \item Les cinq piliers (section~\ref{sec:c30-piliers}).
    \item La méthode~: science du design et falsifiabilité (section~\ref{sec:c30-methode}).
\end{enumerate}

\begin{boiteRemarque}{Ce que ce chapitre est, et n'est pas}
Il expose un programme de recherche déposé en prépublication, non un
savoir établi. Les axiomes et les piliers forment un cadre~; les
prédictions qui en découlent restent à tester. Le
Chapitre~\ref{chap:incertitude} distingue ce qui est établi par des travaux
publiés, ce qui est déposé en prépublication, et ce qui reste à valider.
\end{boiteRemarque}

\separateur

% ============================================================================
\section{L'Objet : la Mission Auditologique}
\label{sec:c30-objet}
% ============================================================================

L'article A0 part d'un constat~: l'audit des SI dispose de méthodes,
de référentiels et d'une littérature instrumentale, mais d'aucun socle
axiomatique ni d'aucun espace de mesure unifié pour son objet central, la
réduction de l'incertitude. Il propose de le reconstituer en
\emph{auditologie}.

\begin{boxdef}
\termecle{Mission auditologique}\index{mission auditologique} (A0)~:
quintuplet
$\mathcal{M} = \langle \mathcal{S}, \mathcal{A}, \mathcal{I},
\mathcal{G}, \mathcal{T} \rangle$, où
\begin{itemize}
    \item $\mathcal{S}$ est le \emph{système audité}, ensemble
          sociotechnique d'agents, d'artefacts, de flux d'information et de
          dispositifs institutionnels~;
    \item $\mathcal{A}$ est l'\emph{architecture cognitive} de l'auditeur
          ou du collectif d'audit, soumise à des contraintes de traitement,
          de mémoire et d'attention~;
    \item $\mathcal{I}$ est le \emph{canal d'information} par lequel les
          preuves sur $\mathcal{S}$ parviennent à $\mathcal{A}$~;
    \item $\mathcal{G}$ est la \emph{structure stratégique} de
          l'interaction entre l'auditeur, l'audité et les tiers~;
    \item $\mathcal{T}$ est la \emph{trajectoire temporelle} du système
          au fil des missions successives.
\end{itemize}

\medskip
\termecle{Auditologie}\index{auditologie}~: science appliquée de la
réduction de l'incertitude sur les missions auditologiques.
\end{boxdef}

\begin{tableaucours}{cL{4.0cm}L{8.0cm}}{%
    \tableauHeader{} &
    \tableauHeader{Composante} &
    \tableauHeader{Dans la mission de LiZbiZ}
}
$\mathcal{S}$ & Système audité & LiZbiZ~: quatre centres de données, douze
agences, l'application intégrée, ses agents et ses procédures \\
$\mathcal{A}$ & Architecture cognitive & L'équipe d'AUDIX CONSEIL, sa
méthode, ses outils O0 à O9 \\
$\mathcal{I}$ & Canal d'information & Entretiens, documents sollicités,
extractions, tests~; ce qui passe et ce qui ne passe pas \\
$\mathcal{G}$ & Structure stratégique & Direction générale, responsable
informatique, magasinier, commissaire aux comptes, auditeurs \\
$\mathcal{T}$ & Trajectoire & Mai~2026, novembre~2026, et les cycles
suivants \\
\end{tableaucours}
\vspace{-0.8em}
\begin{center}{\small\itshape Tableau~30.1. La mission de LiZbiZ comme
mission auditologique}\end{center}

% ============================================================================
\section{Six Axiomes Ontologiques}
\label{sec:c30-axiomes}
% ============================================================================

L'article A0 pose six axiomes qui contraignent toute mission
auditologique. Leur numérotation est celle de l'article, et le cours l'a
suivie depuis le Chapitre~\ref{chap:audit-auditologie}. Les énoncés
ci-dessous reprennent ceux d'A0.

\begin{axiome}[Observabilité partielle]
Pour tout système audité $\mathcal{S}$ non dégénéré, l'auditeur n'a accès
qu'à un sous-ensemble strict des états pertinents de $\mathcal{S}$, et ce
sous-ensemble est lui-même partiellement opaque.
\end{axiome}

\begin{axiome}[Cognition limitée des agents]
L'auditeur et l'audité sont des agents cognitivement limités, au sens de
Simon, de Miller et de Kahneman~\cite{simon1957,miller1956,kahneman2011}.
\end{axiome}

\begin{axiome}[Mesurabilité de l'incertitude épistémique]
L'incertitude épistémique de $\mathcal{A}$ concernant des propositions
spécifiques sur $\mathcal{S}$ admet une mesure quantitative, définissable
par l'entropie informationnelle et ses généralisations.
\end{axiome}

\begin{axiome}[Interaction stratégique]
L'auditeur et l'audité ont des fonctions d'utilité distinctes, et leurs
actions sont, en général, mieux modélisées comme des mouvements dans un jeu
à information imparfaite et asymétrique.
\end{axiome}

\begin{axiome}[Individuation temporelle]
Le système audité $\mathcal{S}$ évolue dans le temps, et sa trajectoire est
dépendante du chemin~: l'état à la mission $n+1$ dépend de l'état à $n$ et
des interventions déclenchées par cette mission antérieure.
\end{axiome}

\begin{axiome}[Opposabilité]
Toute conclusion produite par une mission auditologique doit, en principe,
pouvoir être défendue devant un tiers disposant de l'autorité légitime pour
la contester.
\end{axiome}

\subsection{Ce qui n'est pas un axiome}

A0 écarte explicitement trois candidats. Le cours les a rencontrés.
\begin{itemize}
    \item \textbf{L'indépendance} n'est pas un axiome, mais un théorème,
          dérivable des axiomes~4 et~6 (Chapitre~\ref{chap:normes}).
    \item \textbf{La conformité à un référentiel} (COBIT, ISO/IEC~27001) est
          un choix de fonction objectif (Chapitre~\ref{chap:referentiels}).
    \item \textbf{La matérialité} est un paramètre dérivé de l'opposabilité.
\end{itemize}

% ============================================================================
\section{Quatre Conséquences, Quatre Pratiques}
\label{sec:c30-propositions}
% ============================================================================

A0 en tire quatre propositions. Chacune fonde une pratique du cours.

\begin{proposition}[Modalité d'ignorance explicite]
Sous les axiomes~1, 2, 3 et~4, tout instrument de réponse d'audit valide
doit inclure une modalité explicite pour «~connaissance insuffisante~»
(\NEK{}), distincte de \OUI{} et de \NON{}.
\end{proposition}

\begin{proposition}[Décomposition hiérarchique]
Sous les axiomes~2 et~6, tout instrument couvrant un espace d'audit non
trivial doit être organisé en décomposition hiérarchique, dont la
profondeur est calibrée sur la mémoire de travail de l'auditeur (borne
$7 \pm 2$).
\end{proposition}

\begin{proposition}[Métriques longitudinales]
Sous l'axiome~5 et l'engagement de falsifiabilité, les conclusions d'une
mission ne peuvent être testées empiriquement que par observation
longitudinale.
\end{proposition}

\begin{proposition}[Réflexivité]
Sous les axiomes~4 et~6, l'acte d'audit est lui-même un constituant de la
trajectoire du système audité, et ne peut être modélisé comme une
observation externe.
\end{proposition}

\begin{tableaucours}{L{3.4cm}L{2.4cm}L{6.4cm}}{%
    \tableauHeader{Proposition} &
    \tableauHeader{Axiomes} &
    \tableauHeader{Dans le cours}
}
Modalité \NEK{} & 1, 2, 3, 4 & Instrument \OUI{}/\NON{}/\NEK{}
(Chapitre~\ref{chap:preuve}), mesures $\mu$, $U$, $H$
(Chapitre~\ref{chap:conclure}) \\
Décomposition & 2, 6 & Gabarit canonique et fiche d'assertion
(Chapitre~\ref{chap:decomposition}) \\
Longitudinal & 5, et falsifiabilité & Divergence entre cycles, second
cycle de LiZbiZ (Chapitres~\ref{chap:conclure}
et~\ref{chap:cas-integrateurs}) \\
Réflexivité & 4, 6 & Le suivi des recommandations modifie le système
(Chapitre~\ref{chap:mission})~; le domaine~E change d'observabilité
quand le prestataire change \\
\end{tableaucours}
\vspace{-0.8em}
\begin{center}{\small\itshape Tableau~30.2. Des axiomes aux pratiques du
cours}\end{center}

% ============================================================================
\section{Les Cinq Piliers}
\label{sec:c30-piliers}
% ============================================================================

Les axiomes disent ce qu'est une mission~; les piliers apportent les
théories qui permettent de la modéliser. A0 en retient cinq, chacun avec
un formalisme et une prédiction falsifiable.

\begin{tableaucours}{L{2.4cm}L{3.0cm}L{3.2cm}L{3.6cm}}{%
    \tableauHeader{Pilier} &
    \tableauHeader{Ancrage} &
    \tableauHeader{Formalisme} &
    \tableauHeader{Prédiction falsifiable}
}
1. Cognition limitée & Miller, Simon, Kahneman~\cite{miller1956,simon1957,kahneman2011}
& $\mathcal{A} = \langle K, W, \tau \rangle$, $|W| \le k \approx 7$ & Les
grilles de plus de neuf dimensions simultanées produisent plus
d'incohérences et d'oublis \\
2. Théorie de l'information & Shannon~; Kullback et
Leibler~\cite{shannon1948,kullback1951} & Entropie $H \le \log_2 3$~;
information mutuelle $I(X;Y)$ & Le diagnostic conformité et entropie
prédit mieux que le seul taux de conformité \\
3. Théorie des jeux & von~Neumann, Nash,
Harsanyi~\cite{vonneumann1944,nash1951,harsanyi1967} & Jeu bayésien
$\langle N, (A_i), (\Theta_i), (u_i), p \rangle$ & Les instruments
compatibles avec les incitations réduisent la dissimulation \\
4. Systèmes dynamiques & Chaînes de Markov~\cite{norris1998}, analyse
topologique des données & Niveaux $m_n \in \{1, \dots, 5\}^D$, matrice
de transition $P$ & Les prévisions par matrice de transition battent la
persistance naïve à 12--24 mois \\
5. Ontologie critique-réaliste & Bhaskar, Simondon,
Latour~\cite{bhaskar1975,simondon1958,latour2005} & Strates du réel, de
l'actuel, de l'empirique & Les cadres socio-matériels obtiennent plus de
recommandations mises en œuvre \\
\end{tableaucours}
\vspace{-0.8em}
\begin{center}{\small\itshape Tableau~30.3. Les cinq piliers de
l'auditologie (d'après A0)}\end{center}

\begin{boxinfo}
\textbf{Ce que le cours a utilisé de chaque pilier.} La borne de sept et la
décomposition (pilier~1)~; $H$, $U$ et la divergence (pilier~2)~; l'indice
$\rho$, l'usage stratégique de \NEK{} et l'indépendance (pilier~3)~; la
lecture en cycles et le risque de régression (pilier~4)~; la place des
personnes, des procédures et des artefacts dans chaque constat, par
exemple le compte partagé du magasin de Garoua (pilier~5).
\end{boxinfo}

Les piliers ne sont pas indépendants. Selon A0, la théorie de
l'information présuppose la cognition limitée, la théorie des jeux
présuppose la mesure informationnelle, les systèmes dynamiques
présupposent des équilibres indexés dans le temps, et l'ontologie
sous-tend les quatre autres. Chaque pilier est nécessaire~; leur
suffisance conjointe est une hypothèse de travail. A0 signale aussi une
limite pratique du pilier~4~: l'analyse topologique devient impraticable
au-delà d'une dizaine de dimensions, d'où des premières applications
limitées à huit dimensions au plus.

% ============================================================================
\section{La Méthode : Science du Design et Falsifiabilité}
\label{sec:c30-methode}
% ============================================================================

\subsection{Une science du design}

L'auditologie produit des artefacts~: instruments, métriques, grilles,
logiciels. A0 l'inscrit dans la recherche en science du design~: on
construit un artefact, on l'évalue dans son environnement (cycle de
pertinence), on l'appuie sur la base de connaissances (cycle de rigueur),
et l'on recommence~\cite{hevner2004design}. La méthode de ce cours est
l'\emph{instanciation opérationnelle, pour le praticien,} de ce programme~:
l'instrument \OUI{}/\NON{}/\NEK{}, le gabarit canonique, le classeur O4 et
les labs de NDAP SUKULU sont des artefacts, et chaque promotion
d'étudiants participe à leur évaluation.

\subsection{Un noyau dur et une ceinture protectrice}

A0 adopte le falsificationnisme sophistiqué de Lakatos~\cite{lakatos1970}.
Les six axiomes et les cinq piliers forment un \emph{noyau dur}, qui n'est
pas directement testé. Les hypothèses auxiliaires, les procédures de
mesure et les prédictions forment la \emph{ceinture protectrice}, qui l'est.

\begin{axiome}[Falsifiabilité comme engagement méthodologique]
Toute affirmation spécifique engendrée par l'auditologie doit être énoncée
sous une forme telle que des résultats empiriques puissent, en principe, la
réfuter. La réfutation d'une assertion de la ceinture protectrice n'emporte
pas la réfutation du noyau dur~; elle déclenche une révision de la
ceinture.
\end{axiome}

\begin{tableaucours}{L{2.8cm}L{4.6cm}L{4.8cm}}{%
    \tableauHeader{Pilier} &
    \tableauHeader{Prédiction testable} &
    \tableauHeader{Condition de réfutation}
}
Cognition limitée & Les grilles au-delà de la borne produisent plus
d'erreurs & Aucune différence significative en expérience contrôlée \\
Information & Conformité et entropie prédisent mieux que la conformité
seule & Validité prédictive égale ou inférieure \\
Jeux & Les instruments compatibles avec les incitations réduisent la
dissimulation & Aucune baisse de l'indice d'usage stratégique \\
Dynamique & La matrice de transition bat la persistance naïve & Aucun
avantage prédictif \\
Réalisme critique & Les cadres socio-matériels font mieux appliquer les
recommandations & Aucune différence observable \\
\end{tableaucours}
\vspace{-0.8em}
\begin{center}{\small\itshape Tableau~30.4. Matrice de falsifiabilité
(d'après A0)}\end{center}

\subsection{Ce que l'auditologie n'est pas}

A0 précise quatre refus~: l'auditologie n'est pas \emph{normative} (elle
ne prescrit pas le «~bon audit~»), ni \emph{critique} (elle n'invalide pas
les cadres existants), ni \emph{un référentiel de plus} (ses piliers sont
théoriques), ni \emph{une idéologie} (ses propositions sont falsifiables).
C'est pourquoi ce cours a enseigné ensemble les normes du métier
(Partie~\ref{part:normatif}) et leur lecture auditologique.

\begin{boxlecture}{LiZbiZ, prédiction et constat}
Le mini-cas d'A0 décrit le domaine~D de LiZbiZ~: entropie d'environ
1,2~bit, décomposition en $5 \times 5$ sous-éléments, usage stratégique de
\NEK{} (réponses différentes en entretien formel et en auto-évaluation),
et, à partir d'une matrice de transition estimée sur trois cycles, une
probabilité de régression de 35~\% dans les dix-huit mois. Six mois plus
tard, l'entropie du domaine est d'environ 0,6~bit.

\medskip
\textbf{Constatation.} Le cours retrouve ces valeurs~: $H \approx 1{,}24$
en mai~2026, $H \approx 0{,}57$ en novembre (Chapitre~\ref{chap:conclure}).

\medskip
\textbf{Analyse} (confiance moyenne). L'amélioration est réelle, mais la
prédiction de régression n'est pas encore testée~: elle le sera par les
cycles suivants. C'est exactement la forme d'une affirmation de la
ceinture protectrice~: datée, chiffrée, réfutable. Le domaine~E, qui a
perdu en observabilité au second cycle (Chapitre~\ref{chap:cas-integrateurs}),
rappelle que la trajectoire ne se lit pas sur un seul domaine.

\medskip
\hfill\lecture{0{,}90}{0{,}05}{0{,}57}
\end{boxlecture}

\begin{boxmetier}{Auditeur SI}
\textbf{Ce que l'auditologie change dès la prochaine mission}
\begin{enumerate}
    \item Admettre \NEK{}, et le documenter (axiomes~1 et~3).
    \item Décomposer avant de questionner (axiome~2).
    \item Se demander, pour chaque réponse, qui avait intérêt à la donner
          (axiome~4).
    \item Prévoir le cycle suivant dès le premier rapport (axiome~5).
    \item Écrire chaque conclusion comme si un tiers allait la contester
          (axiome~6).
\end{enumerate}
\end{boxmetier}

% ============================================================================

\separateur

\begin{boxresume}
\textbf{Ce qu'il faut retenir de ce chapitre~:}
\begin{itemize}
    \item \textbf{La mission auditologique} est un quintuplet
          $\langle \mathcal{S}, \mathcal{A}, \mathcal{I}, \mathcal{G},
          \mathcal{T} \rangle$~; l'auditologie est la science appliquée de
          la réduction de l'incertitude sur ces missions.
    \item \textbf{Six axiomes ontologiques}~: observabilité partielle,
          cognition limitée, mesurabilité, interaction stratégique,
          individuation temporelle, opposabilité~; un septième,
          méthodologique~: la falsifiabilité.
    \item \textbf{Ne sont pas des axiomes}~: l'indépendance (un théorème),
          la conformité (un choix d'objectif), la matérialité (un
          paramètre).
    \item \textbf{Quatre propositions} fondent \NEK{}, la décomposition, les
          mesures longitudinales et la réflexivité.
    \item \textbf{Cinq piliers}~: cognition limitée, information, jeux,
          systèmes dynamiques, ontologie critique-réaliste~; chacun avec une
          prédiction réfutable.
    \item \textbf{La méthode du cours} est une instanciation, pour le
          praticien, d'un programme de science du design.
\end{itemize}
\end{boxresume}

\begin{boxinfo}
\textbf{Points de contrôle de l'élève.} Je sais~: énoncer les six axiomes
et l'axiome de falsifiabilité~; dire pourquoi l'indépendance, la
conformité et la matérialité n'en sont pas~; relier chaque proposition à
une pratique du cours~; nommer les cinq piliers et leur prédiction~;
distinguer noyau dur et ceinture protectrice.
\end{boxinfo}

\bigskip

\begin{boxexo}[title={\ding{45}\quad Exercice~30.1. De la pratique à l'axiome \niveaubase}]
Indiquez l'axiome qui fonde chaque pratique.
\begin{enumerate}[(a)]
    \item Coter \NEK{} un document non remis dont on ignore l'existence.
    \item Tenir le contradictoire et conserver la réponse de l'audité.
    \item Comparer une auto-évaluation et un entretien formel.
    \item Prévoir un audit de suivi avec les mêmes questions.
    \item Limiter à sept le nombre d'enfants d'un élément de la
          décomposition.
\end{enumerate}
\end{boxexo}

\begin{boxexo}[title={\ding{45}\quad Exercice~30.2. Démontrer une proposition \niveaumoyen}]
\begin{enumerate}[(a)]
    \item À l'aide des axiomes~1 à~4, expliquez pourquoi une grille
          \OUI{}/\NON{} produit un biais du choix forcé, et pourquoi
          l'indice $\rho$ le mesure (Chapitre~\ref{chap:preuve}).
    \item Expliquez pourquoi l'indépendance se déduit des axiomes~4 et~6.
    \item Pourquoi la conformité à un référentiel ne peut-elle pas être un
          axiome~? Appuyez-vous sur la lecture « Domaine D.3, deux
          critères » du Chapitre~\ref{chap:referentiels}.
\end{enumerate}
\end{boxexo}

\begin{boxexo}[title={\ding{45}\quad Exercice~30.3. Formuler une prédiction réfutable \niveauexpert}]
Vous disposez des deux cycles de LiZbiZ (Chapitre~\ref{chap:cas-integrateurs})
et pouvez en mener deux autres.
\begin{enumerate}[(a)]
    \item Formulez, pour le domaine~E, une prédiction datée et chiffrée,
          rattachée à un pilier.
    \item Précisez la condition de réfutation.
    \item Décrivez le protocole~: questions, critères, calendrier, données
          à conserver.
    \item Si la prédiction est réfutée, qu'est-ce qui est révisé~: le noyau
          dur ou la ceinture protectrice~? Justifiez.
\end{enumerate}
\end{boxexo}

\bigskip

\begin{boxexo}[title={\ding{45}\quad Exercice~30.4. Ce qui n'est pas un axiome \niveaubase}]
Pour l'indépendance, la conformité à un référentiel et la matérialité,
dites ce qu'elles sont selon A0 (théorème, choix d'objectif, paramètre) et
d'où elles se déduisent.
\end{boxexo}

\begin{boxexo}[title={\ding{45}\quad Exercice~30.5. Reconnaître un pilier \niveaumoyen}]
Rattachez à un pilier chaque observation de LiZbiZ~: (a) la grille du
domaine~E est réduite à sept sous-domaines~; (b) le domaine~E passe de
\NEK{} à \NON{} entre deux cycles~; (c) le prestataire répond \NEK{} sur
les questions à fort enjeu~; (d) le compte partagé du magasin de Garoua
mêle une personne, une procédure et un artefact~; (e) le diagnostic utilise
$H$ en plus de $\mu$.
\end{boxexo}

\begin{boxexo}[title={\ding{45}\quad Exercice~30.6. Noyau dur et ceinture \niveaumoyen}]
Une étude montre que des grilles de douze dimensions ne produisent pas plus
d'incohérences que des grilles de sept. Qu'est-ce qui est réfuté~? Que
révise-t-on, et que garde-t-on~? Appuyez-vous sur Lakatos.
\end{boxexo}

\begin{boxexo}[title={\ding{45}\quad Exercice~30.7. Un artefact de science du design \niveaumoyen}]
Décrivez le classeur O4 comme un artefact au sens de la science du
design~: son environnement (cycle de pertinence), sa base de connaissances
(cycle de rigueur), et son évaluation (cycle de conception).
\end{boxexo}

\begin{boxexo}[title={\ding{45}\quad Exercice~30.8. La réflexivité \niveaumoyen}]
Montrez, sur le domaine~E de LiZbiZ, que l'audit du premier cycle a modifié
le système audité. Quelle conséquence pour l'interprétation de la
divergence entre les deux cycles~?
\end{boxexo}

\begin{boxexo}[title={\ding{45}\quad Exercice~30.9. Est-ce réfutable~? \niveauexpert}]
Un collègue soutient que l'auditologie n'est pas une science, parce
qu'aucun résultat ne pourrait réfuter ses axiomes. Répondez en distinguant
la position de Popper et celle de Lakatos, et en citant une prédiction
de la ceinture protectrice.
\end{boxexo}

\begin{boxexo}[title={\ding{45}\quad Exercice~30.10. Ce que l'auditologie ajoute \niveauexpert}]
Comparez les normes professionnelles (Chapitre~\ref{chap:normes}) et
l'auditologie~: que disent les premières, que cherche la seconde~?
Donnez deux pratiques du cours que les normes n'imposent pas et que
l'auditologie justifie, et une limite de l'auditologie.
\end{boxexo}

\bigskip

\begin{boxinfo}
\textbf{Pour aller plus loin.}
\begin{itemize}
    \item L'article A0~\cite{minka2026auditologie}.
    \item I.~Lakatos sur les programmes de recherche~\cite{lakatos1970}.
    \item A.~Hevner et al. sur la science du design~\cite{hevner2004design}.
    \item H.~Simon sur la rationalité limitée~\cite{simon1957}.
\end{itemize}
\end{boxinfo}

\begin{boxinfo}
\textbf{Transition.} Les axiomes demandent de mesurer~; les piliers disent
avec quoi. Le Chapitre~\ref{chap:incertitude} approfondit l'instrument
trichotomique et ses métriques, puis le relie à des travaux publiés sur la
stratégie et sur l'expression de l'incertitude de mesure.
\end{boxinfo}
